\section{Results}
\label{sec:results}

% Generated by scripts/export_results_table.py -- do not edit by hand.
\begin{table}[t]
  \centering
  \caption{Point-prediction quality on the held-out test partition. Hyperparameters for every model were selected on the validation partition; the test partition was read once.}
  \label{tab:results}
  \begin{tabular}{lrrrr}
    \toprule
    Model & Accuracy & MAE (ordinal) & Acc. within 1 & Macro F1 \\
    \midrule
    OrdinalLR & 0.6509 & \textbf{0.3535} & 0.9956 & 0.5635 \\
    Majority & 0.5603 & 0.4397 & 1.0000 & 0.2394 \\
    MultinomialLR & 0.4258 & 0.6298 & 0.9444 & 0.4283 \\
    \bottomrule
  \end{tabular}
\end{table}


Table~\ref{tab:results} reports point-prediction quality on the held-out test
partition. The ordinal model attains the lowest mean absolute ordinal error
by a wide margin over both baselines, confirming that the proportional-odds
structure is doing real work: the multinomial baseline sees identical
features but, lacking the shared-weight ordinal constraint, distributes its
probability mass less efficiently across adjacent levels and pays for it
on every ordinal-aware metric. The majority baseline's accuracy-within-one
score of $1.0000$ is a direct consequence of always predicting "fit" (the
dataset's modal class) under a tolerance of one level either side --
illustrating exactly why within-one accuracy must be read alongside mean
absolute error and macro F1, not instead of them: its macro F1 of
$0.2394$ is the lowest of the three models, reflecting that it is correct
only on the majority class.

% Generated by scripts/export_results_table.py -- do not edit by hand.
\begin{table}[t]
  \centering
  \caption{Split conformal calibration at target coverage $1-\alpha=0.9$, calibrated on the validation partition and evaluated on test.}
  \label{tab:conformal}
  \begin{tabular}{lrrr}
    \toprule
    Model & Target cov. & Empirical cov. & Avg. set size \\
    \midrule
    Majority & 0.9000 & 1.0000 & 3.0000 \\
    MultinomialLR & 0.9000 & 0.8922 & 2.0784 \\
    OrdinalLR & 0.9000 & 0.8855 & 1.7154 \\
    \bottomrule
  \end{tabular}
\end{table}


Table~\ref{tab:conformal} reports the split conformal calibration at a
$90\%$ target coverage, calibrated on the validation partition and
evaluated on test. The ordinal model's conformal sets attain empirical
coverage close to the $90\%$ target while remaining the most informative:
an average set size of $1.72$ out of $3$ possible classes, versus $2.08$
for the multinomial baseline and a trivial $3.0$ (every class, always) for
the majority baseline. The majority baseline's set size directly reflects
that its point probabilities are too flat (identical for every input) to
ever confidently exclude a class -- it "achieves" coverage only by refusing
to narrow the answer at all, the degenerate case the average-set-size metric
exists to catch.

\subsection{Reading the headline number}

The practical takeaway is not "$65\%$ accuracy," which undersells what the
pipeline actually produces. It is: for a held-out shopper the system never
saw during training, the calibrated output is typically a set like
$\{\text{fit}\}$ or $\{\text{fit}, \text{large}\}$ with a guarantee that the
true outcome falls inside that set at least $90\%$ of the time -- a result a
single point prediction cannot make and cannot be checked against.
